
\documentclass[11pt]{article}

\usepackage[margin=1in]{geometry}
\usepackage{amsmath,amssymb,amsthm}
\usepackage{mathtools}
\usepackage{enumitem}
\usepackage{hyperref}

\title{Borel Measurable Sets\\A Summary of Definitions and Properties}
\author{}
\date{}

\begin{document}
	
	\maketitle
	\tableofcontents
	
	\section{Introduction}
	
	Borel measurable sets are among the most fundamental objects in modern
	probability theory, real analysis, and measure theory.
	
	They provide the collection of subsets on which probability measures and
	Lebesgue measure are first defined. Every probability density function,
	distribution function, and measurable random variable is ultimately built
	upon the Borel $\sigma$-algebra.
	
	Throughout this document, we work on the Euclidean space
	
	\[
	\mathbb{R}^n.
	\]
	
	%%%%%%%%%%%%%%%%%%%%%%%%%%%%%%%%%%%%%%%%%%%%%%%%%%%%%%%%%%%%%%
	
	\section{Open and Closed Sets}
	
	Everything begins with the notion of an open set.
	
	\subsection{Open Sets}
	
	A subset
	
	\[
	U\subseteq\mathbb R^n
	\]
	
	is called \textbf{open} if for every point
	
	\[
	x\in U,
	\]
	
	there exists some radius
	
	\[
	r>0
	\]
	
	such that
	
	\[
	B_r(x)
	=
	\{y\in\mathbb R^n:\|x-y\|<r\}
	\subseteq U.
	\]
	
	In words, every point has a small neighborhood completely contained in the
	set.
	
	Examples include
	
	\[
	(a,b),
	\qquad
	(-\infty,c),
	\qquad
	\mathbb R^n.
	\]
	
	\subsection{Closed Sets}
	
	A subset
	
	\[
	F\subseteq\mathbb R^n
	\]
	
	is called \textbf{closed} if its complement is open.
	
	Examples include
	
	\[
	[a,b],
	\qquad
	[c,\infty),
	\qquad
	\emptyset.
	\]
	
	%%%%%%%%%%%%%%%%%%%%%%%%%%%%%%%%%%%%%%%%%%%%%%%%%%%%%%%%%%%%%%
	
	\section{The Borel $\sigma$-Algebra}
	
	\subsection{Definition}
	
	The collection of all Borel measurable sets is called the
	\textbf{Borel $\sigma$-algebra}.
	
	It is denoted by
	
	\[
	\boxed{
		\mathcal B(\mathbb R^n).
	}
	\]
	
	It is defined as
	
	\[
	\boxed{
		\mathcal B(\mathbb R^n)
		=
		\sigma(\text{all open sets}).
	}
	\]
	
	This means that it is the smallest $\sigma$-algebra containing every open
	subset of $\mathbb R^n$.
	
	%%%%%%%%%%%%%%%%%%%%%%%%%%%%%%%%%%%%%%%%%%%%%%%%%%%%%%%%%%%%%%
	
	\section{What is a $\sigma$-Algebra?}
	
	A collection
	
	\[
	\mathcal F
	\]
	
	of subsets of a set
	
	\[
	\Omega
	\]
	
	is called a \textbf{$\sigma$-algebra} if it satisfies the following three
	properties.
	
	\subsection*{1. Contains the Whole Space}
	
	\[
	\boxed{
		\Omega\in\mathcal F.
	}
	\]
	
	Consequently,
	
	\[
	\emptyset\in\mathcal F.
	\]
	
	\subsection*{2. Closed Under Complements}
	
	Whenever
	
	\[
	A\in\mathcal F,
	\]
	
	then
	
	\[
	\boxed{
		A^c\in\mathcal F.
	}
	\]
	
	\subsection*{3. Closed Under Countable Unions}
	
	If
	
	\[
	A_1,A_2,\ldots\in\mathcal F,
	\]
	
	then
	
	\[
	\boxed{
		\bigcup_{n=1}^{\infty}A_n
		\in
		\mathcal F.
	}
	\]
	
	%%%%%%%%%%%%%%%%%%%%%%%%%%%%%%%%%%%%%%%%%%%%%%%%%%%%%%%%%%%%%%
	
	\section{Immediate Consequences}
	
	From the axioms above one can prove that every $\sigma$-algebra is also
	closed under
	
	\subsection*{Countable Intersections}
	
	Using De Morgan's Laws,
	
	\[
	\boxed{
		\bigcap_{n=1}^{\infty}A_n
		\in
		\mathcal F.
	}
	\]
	
	\subsection*{Finite Unions}
	
	\[
	A\cup B\in\mathcal F.
	\]
	
	\subsection*{Finite Intersections}
	
	\[
	A\cap B\in\mathcal F.
	\]
	
	\subsection*{Set Differences}
	
	\[
	\boxed{
		A\setminus B
		=
		A\cap B^c
		\in
		\mathcal F.
	}
	\]
	
	%%%%%%%%%%%%%%%%%%%%%%%%%%%%%%%%%%%%%%%%%%%%%%%%%%%%%%%%%%%%%%
	
	\section{Equivalent Generators}
	
	Although the Borel $\sigma$-algebra is defined using open sets, it can be
	generated from many different collections.
	
	Each of the following generates exactly the same $\sigma$-algebra:
	
	\[
	\boxed{
		\mathcal B(\mathbb R)
		=
		\sigma(\text{open sets})
	}
	\]
	
	\[
	=
	\sigma(\text{closed sets})
	\]
	
	\[
	=
	\sigma(\text{open intervals})
	\]
	
	\[
	=
	\sigma(\text{closed intervals})
	\]
	
	\[
	=
	\sigma((-\infty,a))
	\]
	
	\[
	=
	\sigma((-\infty,a])
	\]
	
	\[
	=
	\sigma((a,\infty))
	\]
	
	\[
	=
	\sigma([a,\infty)).
	\]
	
	Thus, any one of these families may be used to define the Borel
	$\sigma$-algebra.
	
	%%%%%%%%%%%%%%%%%%%%%%%%%%%%%%%%%%%%%%%%%%%%%%%%%%%%%%%%%%%%%%
	
	\section{Examples of Borel Sets}
	
	Examples include
	
	\[
	(a,b),
	\]
	
	\[
	[a,b],
	\]
	
	\[
	(a,b],
	\]
	
	\[
	[a,b),
	\]
	
	finite unions of intervals,
	
	countable unions,
	
	countable intersections,
	
	singletons
	
	\[
	\{x\},
	\]
	
	finite sets,
	
	countable sets,
	
	the integers
	
	\[
	\mathbb Z,
	\]
	
	the rational numbers
	
	\[
	\mathbb Q,
	\]
	
	the irrational numbers
	
	\[
	\mathbb R\setminus\mathbb Q,
	\]
	
	and every open or closed subset of
	
	\[
	\mathbb R^n.
	\]
	
	%%%%%%%%%%%%%%%%%%%%%%%%%%%%%%%%%%%%%%%%%%%%%%%%%%%%%%%%%%%%%%
	
	\section{Examples of Non-Borel Sets}
	
	The Borel $\sigma$-algebra is extremely large; it contains every open set,
	every closed set, every interval, every countable set, and every set that can
	be obtained from these through countable unions, countable intersections, and
	complements.
	
	Nevertheless,
	
	\[
	\boxed{
		\mathcal B(\mathbb R)
		\neq
		\mathcal P(\mathbb R),
	}
	\]
	
	where
	
	\[
	\mathcal P(\mathbb R)
	\]
	
	denotes the power set of the real numbers.
	
	In other words, there exist subsets of the real line that are not Borel
	measurable.
	
	These sets cannot be explicitly written down by a finite formula; their
	existence relies on the Axiom of Choice.
	
	We now discuss three classical examples.
	
	%%%%%%%%%%%%%%%%%%%%%%%%%%%%%%%%%%%%%%%%%%%%%%%%%%%%%%%%%%%%%%
	
	\subsection{Vitali Sets}
	
	Perhaps the most famous example is the \emph{Vitali set}.
	
	\paragraph{Step 1. Define an Equivalence Relation}
	
	Define a relation on $\mathbb R$ by
	
	\[
	x\sim y
	\quad
	\Longleftrightarrow
	\quad
	x-y\in\mathbb Q.
	\]
	
	This is an equivalence relation because
	
	\begin{enumerate}
		\item
		$x-x=0\in\mathbb Q$,
		
		\item
		if $x-y\in\mathbb Q$, then
		
		\[
		y-x=-(x-y)\in\mathbb Q,
		\]
		
		\item
		if
		
		\[
		x-y\in\mathbb Q,
		\qquad
		y-z\in\mathbb Q,
		\]
		
		then
		
		\[
		x-z=(x-y)+(y-z)\in\mathbb Q.
		\]
	\end{enumerate}
	
	Thus,
	
	\[
	\mathbb R
	\]
	
	is partitioned into equivalence classes.
	
	%%%%%%%%%%%%%%%%%%%%%%%%%%%%%%%%%%%%%%%%%%%%%%%%%%%%%%%%%%%%%%
	
	\paragraph{Step 2. Choose One Representative}
	
	Using the Axiom of Choice, choose exactly one point from every equivalence
	class.
	
	The resulting set is called a Vitali set,
	
	\[
	V.
	\]
	
	%%%%%%%%%%%%%%%%%%%%%%%%%%%%%%%%%%%%%%%%%%%%%%%%%%%%%%%%%%%%%%
	
	\paragraph{Step 3. Main Properties}
	
	Every real number differs from exactly one element of
	
	\[
	V
	\]
	
	by a rational number.
	
	\paragraph{Why Every Real Number Differs from Exactly One Element of \(V\) by a Rational Number}
	
	Recall the equivalence relation
	
	\[
	x\sim y
	\quad\Longleftrightarrow\quad
	x-y\in\mathbb Q.
	\]
	
	This relation partitions the real line into disjoint equivalence classes.
	
	For every real number \(x\), its equivalence class is
	
	\[
	[x]
	=
	\{\,x+q:q\in\mathbb Q\,\}.
	\]
	
	Thus, two real numbers belong to the same equivalence class if and only if
	their difference is rational.
	
	\bigskip
	
	Using the Axiom of Choice, we construct a set
	
	\[
	V\subseteq\mathbb R
	\]
	
	by selecting exactly one representative from each equivalence class.
	
	That is,
	
	\[
	\boxed{
		\text{For every equivalence class }[x],
		\text{ there exists exactly one }
		v\in V
		\text{ such that }
		v\in[x].
	}
	\]
	
	This single property immediately implies the following important result.
	
	\bigskip
	
	\paragraph{Existence}
	
	Let
	
	\[
	x\in\mathbb R
	\]
	
	be arbitrary.
	
	Since the equivalence classes partition the real numbers,
	
	\[
	x
	\]
	
	belongs to exactly one equivalence class,
	
	\[
	[x].
	\]
	
	By construction,
	
	\[
	V
	\]
	
	contains one representative of this class.
	
	Call this representative
	
	\[
	v.
	\]
	
	Because both
	
	\[
	x
	\quad\text{and}\quad
	v
	\]
	
	belong to the same equivalence class,
	
	\[
	x\sim v.
	\]
	
	Therefore,
	
	\[
	x-v\in\mathbb Q.
	\]
	
	Equivalently,
	
	\[
	\boxed{
		x=v+q
	}
	\]
	
	for some rational number
	
	\[
	q\in\mathbb Q.
	\]
	
	Hence every real number differs from at least one element of \(V\) by a
	rational number.
	
	\bigskip
	
	\paragraph{Uniqueness}
	
	Suppose that there exist two elements
	
	\[
	v_1,v_2\in V
	\]
	
	such that
	
	\[
	x-v_1\in\mathbb Q,
	\qquad
	x-v_2\in\mathbb Q.
	\]
	
	Subtracting the two relations gives
	
	\[
	v_1-v_2
	=
	(v_1-x)+(x-v_2).
	\]
	
	Since the rational numbers are closed under addition,
	
	\[
	v_1-v_2\in\mathbb Q.
	\]
	
	Therefore,
	
	\[
	v_1\sim v_2.
	\]
	
	Thus,
	
	\[
	v_1
	\quad\text{and}\quad
	v_2
	\]
	
	belong to the same equivalence class.
	
	However, the set
	
	\[
	V
	\]
	
	was constructed to contain exactly one representative from each equivalence
	class.
	
	Consequently,
	
	\[
	\boxed{
		v_1=v_2.
	}
	\]
	
	This proves that the representative is unique.
	
	\bigskip
	
	Combining existence and uniqueness, we conclude that
	
	\[
	\boxed{
		\forall x\in\mathbb R,\;
		\exists!\,v\in V
		\text{ such that }
		x-v\in\mathbb Q.
	}
	\]
	
	The symbol
	
	\[
	\exists!
	\]
	
	means "there exists exactly one."
	
	Equivalently, every real number can be written uniquely in the form
	
	\[
	\boxed{
		x=v+q,
		\qquad
		v\in V,\;
		q\in\mathbb Q.
	}
	\]
	
	where \(v\) is the unique representative of the equivalence class containing
	\(x\).
	
	Consequently,
	
	\[
	\mathbb R
	=
	\bigcup_{q\in\mathbb Q}
	(V+q),
	\]
	
	where
	
	\[
	V+q
	=
	\{v+q:v\in V\}.
	\]
	
	Moreover,
	
	\[
	(V+q_1)\cap(V+q_2)=\varnothing,
	\qquad
	q_1\neq q_2.
	\]
	
	%%%%%%%%%%%%%%%%%%%%%%%%%%%%%%%%%%%%%%%%%%%%%%%%%%%%%%%%%%%%%%
	
	\paragraph{Step 4. Why It Is Not Borel}
	
	Every Borel set is Lebesgue measurable.
	
	However, one proves that no Vitali set is Lebesgue measurable.
	
	Therefore,
	
	\[
	\boxed{
		V
		\notin
		\mathcal B(\mathbb R).
	}
	\]
	
	%%%%%%%%%%%%%%%%%%%%%%%%%%%%%%%%%%%%%%%%%%%%%%%%%%%%%%%%%%%%%%
	
	\subsection{Bernstein Sets}
	
	A second important example is the Bernstein set.
	
	A subset
	
	\[
	B\subseteq\mathbb R
	\]
	
	is called a Bernstein set if
	
	\begin{enumerate}
		\item
		every nonempty perfect set intersects \(B\),
		
		\item
		every nonempty perfect set also intersects
		
		\[
		\mathbb R\setminus B.
		\]
	\end{enumerate}
	
	Thus,
	
	\[
	B
	\]
	
	contains no entire perfect set, but neither does its complement.
	
	%%%%%%%%%%%%%%%%%%%%%%%%%%%%%%%%%%%%%%%%%%%%%%%%%%%%%%%%%%%%%%
	
	\paragraph{Existence}
	
	Bernstein sets exist only because of the Axiom of Choice.
	
	Their construction requires a transfinite induction over all perfect subsets
	of the real line.
	
	%%%%%%%%%%%%%%%%%%%%%%%%%%%%%%%%%%%%%%%%%%%%%%%%%%%%%%%%%%%%%%
	
	\paragraph{Why They Are Not Borel}
	
	A fundamental theorem of descriptive set theory states:
	
	\medskip
	
	\emph{Every uncountable Borel subset of $\mathbb R$ contains a nonempty
		perfect subset.}
	
	\medskip
	
	If a Bernstein set were Borel, then it would contain a perfect subset.
	
	This contradicts its defining property.
	
	Hence,
	
	\[
	\boxed{
		B
		\notin
		\mathcal B(\mathbb R).
	}
	\]
	
	%%%%%%%%%%%%%%%%%%%%%%%%%%%%%%%%%%%%%%%%%%%%%%%%%%%%%%%%%%%%%%
	
	\subsection{Hamel Bases}
	
	A Hamel basis is a basis of the vector space
	
	\[
	\mathbb R
	\]
	
	over the field
	
	\[
	\mathbb Q.
	\]
	
	Thus,
	
	every real number admits a unique finite representation
	
	\[
	x
	=
	q_1b_1+\cdots+q_nb_n,
	\]
	
	where
	
	\[
	q_i\in\mathbb Q
	\]
	
	and
	
	\[
	b_i
	\]
	
	belong to the Hamel basis.
	
	%%%%%%%%%%%%%%%%%%%%%%%%%%%%%%%%%%%%%%%%%%%%%%%%%%%%%%%%%%%%%%
	
	\paragraph{Existence}
	
	Every vector space possesses a basis.
	
	The proof uses Zorn's Lemma, which is equivalent to the Axiom of Choice.
	
	%%%%%%%%%%%%%%%%%%%%%%%%%%%%%%%%%%%%%%%%%%%%%%%%%%%%%%%%%%%%%%
	
	\paragraph{Why Hamel Bases Are Not Borel}
	
	A deep theorem states that every Hamel basis is non-Lebesgue measurable.
	
	Since every Borel set is Lebesgue measurable, it follows immediately that
	
	\[
	\boxed{
		\text{Every Hamel basis is not Borel measurable.}
	}
	\]
	
	The proof relies on the Steinhaus Theorem and the algebraic structure of
	Hamel bases.
	
	%%%%%%%%%%%%%%%%%%%%%%%%%%%%%%%%%%%%%%%%%%%%%%%%%%%%%%%%%%%%%%
	
	\subsection{Summary}
	
	These examples illustrate an important fact.
	
	Although the Borel $\sigma$-algebra contains virtually every set encountered
	in analysis, it is still strictly smaller than the collection of all subsets
	of the real numbers.
	
	The existence of non-Borel sets depends essentially on the Axiom of Choice
	and cannot be demonstrated constructively.
	
	The three classical examples are summarized below.
	
	\[
	\begin{array}{|c|c|c|}
		\hline
		\textbf{Set} &
		\textbf{Exists by} &
		\textbf{Reason it is not Borel}
		\\
		\hline
		\text{Vitali set}
		&
		\text{Axiom of Choice}
		&
		\text{Not Lebesgue measurable}
		\\
		\hline
		\text{Bernstein set}
		&
		\text{Transfinite induction}
		&
		\text{Contains no perfect subset}
		\\
		\hline
		\text{Hamel basis}
		&
		\text{Zorn's Lemma}
		&
		\text{Not Lebesgue measurable}
		\\
		\hline
	\end{array}
	\]
	
	%%%%%%%%%%%%%%%%%%%%%%%%%%%%%%%%%%%%%%%%%%%%%%%%%%%%%%%%%%%%%%
	
	\section{Why Are Borel Sets Important?}
	
	Probability measures are defined on Borel sets.
	
	If
	
	\[
	X
	\]
	
	is a random variable, then for every Borel set
	
	\[
	B,
	\]
	
	the probability
	
	\[
	P(X\in B)
	\]
	
	is well defined.
	
	Likewise,
	
	Lebesgue measure,
	
	Gaussian measures,
	
	Poisson processes,
	
	Brownian motion,
	
	and nearly every object in probability theory are defined on Borel
	$\sigma$-algebras.
	
	%%%%%%%%%%%%%%%%%%%%%%%%%%%%%%%%%%%%%%%%%%%%%%%%%%%%%%%%%%%%%%
	
	\section{Borel Measurable Functions}
	
	A function
	
	\[
	f:\mathbb R^n\to\mathbb R
	\]
	
	is called \textbf{Borel measurable} if
	
	\[
	\boxed{
		f^{-1}(B)
		\in
		\mathcal B(\mathbb R^n)
	}
	\]
	
	for every Borel set
	
	\[
	B
	\subseteq
	\mathbb R.
	\]
	
	Equivalently,
	
	\[
	f^{-1}((-\infty,a))
	\]
	
	must be a Borel set for every real number
	
	\[
	a.
	\]
	
	Every continuous function is Borel measurable.
	
	%%%%%%%%%%%%%%%%%%%%%%%%%%%%%%%%%%%%%%%%%%%%%%%%%%%%%%%%%%%%%%
	
	\section{Summary of Properties}
	
	\begin{center}
		\renewcommand{\arraystretch}{1.5}
		\begin{tabular}{|l|l|}
			\hline
			Property & Holds?\\
			\hline
			Contains every open set & Yes\\
			Contains every closed set & Yes\\
			Closed under complements & Yes\\
			Closed under countable unions & Yes\\
			Closed under countable intersections & Yes\\
			Closed under finite unions & Yes\\
			Closed under finite intersections & Yes\\
			Contains every interval & Yes\\
			Contains every singleton & Yes\\
			Contains every countable set & Yes\\
			Contains every continuous preimage & Yes\\
			Contains every subset of $\mathbb R$ & No\\
			\hline
		\end{tabular}
	\end{center}
	
	%%%%%%%%%%%%%%%%%%%%%%%%%%%%%%%%%%%%%%%%%%%%%%%%%%%%%%%%%%%%%%
	
	\section{Key Facts to Remember}
	
	\begin{enumerate}[label=\arabic*.]
		\item The Borel $\sigma$-algebra is generated by the open sets.
		
		\item It is the smallest $\sigma$-algebra containing all open sets.
		
		\item Every open, closed, and interval-type set is Borel measurable.
		
		\item Every continuous function is Borel measurable.
		
		\item Every probability measure on $\mathbb R^n$ is defined on Borel sets.
		
		\item Random variables are measurable precisely so that probabilities of the
		form
		
		\[
		P(X\in B)
		\]
		
		are well defined for every Borel set
		
		\[
		B.
		\]
	\end{enumerate}
	
\end{document}

